% Template for DAN, KLTN, SVNCKH 

\documentclass[oneside,a4paper,13pt]{extreport}

% Font tiếng Việt
\usepackage[T5]{fontenc}
\usepackage[utf8]{inputenc}
\DeclareTextSymbolDefault{\DH}{T1}

% Tài liệu tham khảo
\usepackage[
	sorting=nty,
	backend=bibtex,
	defernumbers=true]{biblatex}
\usepackage[unicode]{hyperref} % Bookmark tiếng Việt
\addbibresource{References/references.bib}

\makeatletter
\def\blx@maxline{77}
\makeatother

% Chèn hình, các hình trong luận văn được để trong thư mục Images/
\usepackage{graphicx}
\graphicspath{ {images/} }

% Chèn và định dạng mã nguồn
\usepackage{listings}
\usepackage{color}
\definecolor{codegreen}{rgb}{0,0.6,0}
\definecolor{codegray}{rgb}{0.5,0.5,0.5}
\definecolor{codepurple}{rgb}{0.58,0,0.82}
\definecolor{backcolour}{rgb}{0.95,0.95,0.92}
\lstdefinestyle{mystyle}{
    backgroundcolor=\color{backcolour},   
    commentstyle=\color{codegreen},
    keywordstyle=\color{magenta},
    numberstyle=\tiny\color{codegray},
    stringstyle=\color{codepurple},
    basicstyle=\footnotesize,
    breakatwhitespace=false,         
    breaklines=true,                 
    captionpos=b,                    
    keepspaces=true,                 
    numbers=left,                    
    numbersep=5pt,                  
    showspaces=false,                
    showstringspaces=false,
    showtabs=false,                  
    tabsize=2
}
\lstset{style=mystyle}

% Chèn và định dạng mã giả
\usepackage{amsmath}
\usepackage{algorithm}
\usepackage[noend]{algpseudocode}
\makeatletter
\def\BState{\State\hskip-\ALG@thistlm}
\makeatother

% Bảng biểu
\usepackage{multirow}
\usepackage{array}
\newcolumntype{L}[1]{>{\raggedright\let\newline\\\arraybackslash\hspace{0pt}}m{#1}}
\newcolumntype{C}[1]{>{\centering\let\newline\\\arraybackslash\hspace{0pt}}m{#1}}
\newcolumntype{R}[1]{>{\raggedleft\let\newline\\\arraybackslash\hspace{0pt}}m{#1}}



 
% Đổi tên mặc định
\renewcommand{\chaptername}{Chương}
\renewcommand{\figurename}{Hình}
\renewcommand{\tablename}{Bảng}
\renewcommand{\contentsname}{Mục lục}
\renewcommand{\listfigurename}{Danh sách hình}
\renewcommand{\listtablename}{Danh sách bảng}
\renewcommand{\appendixname}{Phụ lục}

% Định dạng chapter
\usepackage{titlesec}
\titleformat{\chapter}
    [display] 
    {\normalfont\bfseries\large}{\chaptername \ \thechapter}{10pt}{\huge}
\titlespacing*{\chapter}{0pt}{-10pt}{20pt} %khoảng cách giữa chapter và đầu trang

\titleformat{\section}
    {\normalfont\bfseries\large}{\thesection}{1em}{}
 
\titleformat{\subsection}
    {\normalfont\bfseries\normalsize}{\thesubsection}{1em}{}
  
% Dãn dòng 1.5
\usepackage{setspace}
\onehalfspacing

% Thụt vào đầu dòng
\usepackage{indentfirst}

% Canh lề
\usepackage[
  top=35mm,
  bottom=20mm,
  left=35mm,
  right=20mm,
  footskip = 15mm,
  includefoot]{geometry}
  
% Trang bìa
\usepackage{tikz}
\usetikzlibrary{calc}
\newcommand\HRule{\rule{\textwidth}{1pt}}

% ========================================================================================= %
% CHÚ Ý: Thông tin chung về KLTN - sinh viên điền vào đây để tự động update các trang khác  %
% ========================================================================================= %
\newcommand{\tenSV}{Nguyễn~Văn~A} % Dấu ~ là khoảng trắng không được tách (các chữ nối với nhau bằng dấu ~ sẽ nằm cùng 1 dòng
\newcommand{\mssv}{1234567}
\newcommand{\tenKL}{Sử~dụng~LaTeX trong Khoá~luận~tốt~nghiệp} % Chú ý dấu ~ trong tên khóa luận
\newcommand{\tenGVHD}{Tên~Giáo~Viên}
\newcommand{\tenBM}{Công nghệ tri thức}

\begin{document}

\begin{titlepage}

\begin{center}
{BỘ GIÁO DỤC VÀ ĐÀO TẠO}\\
{TRƯỜNG ĐẠI HỌC MỞ THÀNH PHỐ HỒ CHÍ MINH}\\[1cm]
%\textbf{KHOA CÔNG NGHỆ THÔNG TIN}\\[2cm]

\includegraphics[scale=0.15]{images/logo}\\[1.5cm]

{ \Large \bfseries NGUYỄN VĂN A\\[2.5cm] } 

%Tên đề tài Khóa luận tốt nghiệp/Đồ án tốt nghiệp
{ \huge \bfseries MẪU BÁO CÁO\\KHÓA LUẬN TỐT NGHIỆP
%,\\ĐỒ ÁN NGÀNH 
\\[2.5cm]} 


%Chọn một trong các dòng sau
\large KHÓA LUẬN TỐT NGHIỆP\\
%\large ĐỒ ÁN NGÀNH\\

%Chọn một trong các dòng sau
\large NGÀNH CÔNG NGHỆ THÔNG TIN%\\[2.cm]
%\large NGÀNH KHOA HỌC MÁY TÍNH\\
%\large NGÀNH HỆ THỐNG THÔNG TIN QUẢN LÝ\\
%\large NGÀNH TRÍ TUỆ NHÂN TẠO\\

\begin{tikzpicture}[remember picture, overlay]
  \draw[line width = 2pt] ($(current page.north west) + (3.5cm,-2cm)$) rectangle ($(current page.south east) + (-1.5cm,3cm)$);
\end{tikzpicture}

\vfill
TP. Hồ Chí Minh, YYYY

\end{center}

\pagebreak


\thispagestyle{empty}
\begin{center}

BỘ GIÁO DỤC VÀ ĐÀO TẠO\\
{TRƯỜNG ĐẠI HỌC MỞ THÀNH PHỐ HỒ CHÍ MINH}\\[1.cm]

\includegraphics[scale=0.15]{images/logo}\\[1.cm]

{\large \bfseries NGUYỄN VĂN A - 123456789\\[2cm]}

%Tên đề tài Khóa luận tốt nghiệp/Đồ án ngành
{ \huge \bfseries MẪU BÁO CÁO\\KHÓA LUẬN TỐT NGHIỆP
%,\\ĐỒ ÁN NGÀNH
\\[2cm] } 


%Chọn một trong các dòng sau
\large KHÓA LUẬN TỐT NGHIỆP\\
%\large ĐỒ ÁN TỐT NGHIỆP CỬ NHÂN\\

%Chọn một trong các dòng sau
\large NGÀNH CÔNG NGHỆ THÔNG TIN\\[1.5cm]
%\large NGÀNH KHOA HỌC MÁY TÍNH\\[2cm]
%\large NGÀNH HỆ THỐNG THÔNG TIN QUẢN LÝ\\[2cm]
%\large NGÀNH TRÍ TUỆ NHÂN TẠO\\[2cm]


Giảng viên hướng dẫn:\\
\textbf{TS. LÊ VIẾT TUẤN}


\begin{tikzpicture}[remember picture, overlay]
  \draw[line width = 2pt] ($(current page.north west) + (3.5cm,-2cm)$) rectangle ($(current page.south east) + (-1.5cm,3cm)$);
\end{tikzpicture}

\vfill
TP. Hồ Chí Minh, YYYY

\end{center}

\end{titlepage}

\pagenumbering{roman} % Đánh số i, ii, iii, ...

\addcontentsline{toc}{chapter}{Ý kiến Giảng viên hướng dẫn}
\chapter*{Ý kiến cho phép bảo vệ ĐỒ Án/ Khóa luận tốt nghiệp của Giảng viên hướng dẫn}
\label{adviser}

Giảng viên hướng dẫn: Lê Viết Tuấn

Sinh viên thực hiện: Nguyễn Văn A
...

\addcontentsline{toc}{chapter}{Lời cảm ơn}
\chapter*{Lời cảm ơn}
\label{thanks}

Tôi xin chân thành cảm ơn ...

% Mục lục, danh sách hình, danh sách bảng
\addcontentsline{toc}{chapter}{Mục lục}
\tableofcontents
\listoffigures
\listoftables

\addcontentsline{toc}{chapter}{Tóm tắt}
\include{Appendix/tomtat}

\clearpage

\pagenumbering{arabic} % Đánh số 1, 2, 3, ...

% Các chương nội dung
\chapter{Giới thiệu}
\label{Chapter1}
Đây là mẫu báo cáo Sinh viên nghiên cứu khoa học, Đồ án ngành, Khóa luận tốt nghiệp dành cho sinh viên Khoa Công nghệ thông tin, Trường Đại học Mở Thành phố Hồ Chí Minh.

Sinh viên khi sử dụng mẫu này, cần tùy chỉnh cho phù hợp với cá nhân mình, phù hợp với quy định của Khoa, của Trường tại thời điểm viết báo cáo.


\chapter{Trình bày báo cáo}
\label{Chapter2}

\section{Quy định chung}
Báo cáo phải được trình bày ngắn gọn, rõ ràng, mạch lạc, sạch sẽ, không được tẩy xóa, có đánh số trang, đánh số bảng biểu, hình vẽ, đồ thị. 

Nội dung báo cáo phải được phân thành các chương. Số thứ tự của các chương, mục được đánh số bằng hệ thống số Ả-rập, không dùng số La mã. Các mục và tiểu mục được đánh số bằng các nhóm hai hoặc ba chữ số, cách nhau một dấu chấm: số thứ nhất chỉ số chương, chỉ số thứ hai chỉ số mục, số thứ ba chỉ số tiểu mục.

Báo cáo trình bày sử dụng khổ giấy với việc canh lề như sau: Lề trên 3,5 cm, lề dưới 3.0 cm, lề trái 3,5 cm, lề phải 2 cm. 

Font chữ dùng trong báo cáo (Times New Roman) với kích cỡ (size) 13pt, sử dụng chế độ dãn dòng (line spacing) chế độ 1.5 lines.

\section{Bố cục của báo cáo}

Một báo cáo thường gồm các chương sau:

\begin{itemize}
\item \textbf{Giới thiệu (Chương 1)}: 

- Giới thiệu tổng quan về đề tài, lý do chọn đề tài, mục đích, đối tượng và phạm vi nghiên cứu.

- Mô tả bài toán mà đề tài giải quyết và những hướng giải quyết hiện có của đề tài.

- Những hạn chế, những vấn đề của những hướng giải quyết trước đây là gì?

- Những vấn đề mà đề tài sẽ giải quyết.

- Các đóng góp chính của đề tài.

\item \textbf{Các công trình liên quan (Chương 2)}: 

- Phân tích đánh giá các đề tài/ nghiên cứu đã có của các tác giả trong và ngoài nước liên quan đến đề tài.

- Từ những vấn đề còn tồn tại, chỉ ra những vấn đề mà đề tài sẽ tập trung, nghiên cứu giải quyết.

\item \textbf{Phương pháp đề xuất (Chương 3)}

- Trình bày cơ sở lý thuyết, lý luận, giả thiết khoa học và phương pháp nghiên cứu đã được sử dụng trong đề tài.

- Nếu đề xuất hướng giải quyết mới, mô hình mới thì cần mô tả chi tiết phương pháp mới đó.

\item \textbf{Kết quả thí nghiệm (Chương 4)}

- Mô tả các kết quả nghiên cứu khoa học hoặc kết quả thực nghiệm. Thông thường cần mô tả môi trường thí nghiệm trước như mô tả các tập dữ liệu để kiểm chứng phương pháp/ mô hình. Miêu tả các độ đo để đánh giá kết quả. Môi trường chạy thí nghiệm. So sánh kết quả thực nghiệm với các nghiên cứu liên quan trước đó, bàn luận và giải thích kết quả.

- Đối với đề tài ứng dụng có kết quả là sản phẩm phần mềm phải có hồ sơ thiết kế, cài đặt,\ldots theo một trong các mô hình đã học (UML,\ldots).

\item \textbf{Kết luận (Chương 5)}:

 - Tóm tắt lại những kết quả đạt được, những đóng góp mới và những đề xuất mới, \ldots.
 
 - Những hạn chế mà đề tài chưa giải quyết được (nếu có) và hướng phát triển trong tương lai.

\item \textbf{Tài liệu tham khảo}: 

- Chỉ bao gồm các tài liệu được trích dẫn, sử dụng và đề cập tới để bàn luận trong báo cáo.

- Hướng dẫn trích dẫn ở chương \ref{Chapter4}.

\item \textbf{Phụ lục (nếu có)}:

- Phần này bao gồm nội dung cần thiết nhằm minh họa hoặc hỗ trợ cho nội dung báo cáo như số liệu, mẫu biểu, tranh ảnh,\ldots 

- Nếu có công trình công bố thì sẽ để vào phần phụ lục này.
\end{itemize}



\chapter{Bảng biểu, hình vẽ, phương trình}
\label{Chapter3}
%\section{Bảng biểu, hình vẽ, phương trình}
Việc đánh số bảng biểu, hình vẽ, phương trình phải gắn với số chương; ví dụ hình 3.4 có nghĩa là hình thứ 4 trong Chương 3.
Mọi đồ thị, bảng biểu, hình vẽ lấy từ các nguồn khác phải được trích dẫn đầy đủ. Đầu đề của bảng biểu ghi \textbf{phía trên} bảng, đầu đề của hình vẽ ghi \textbf{phía dưới} hình. 
Khi đề cập đến các bảng biểu và hình vẽ phải nêu rõ số của hình và bảng biểu đó, ví dụ ``... được nêu trong Bảng 4.1'' hoặc ``xem Hình 3.2'' mà không được viết ``… được nêu trong bảng dưới đây'' hoặc ``trong đồ thị của X và Y sau''.

\section{Bảng biểu}
Thông thường, những bảng ngắn và đồ thị phải đi liền với phần nội dung đề cập tới các bảng và đồ thị này ở lần thứ nhất.
Các bảng dài có thể để ở những trang riêng nhưng cũng phải tiếp theo ngay phần nội dung đề cập tới bảng này ở lần đầu tiên.

\begin{table}[ht]
\caption{Ví dụ về bảng biểu 
\label{tab:table1}}
\centering
\begin{tabular}{|c||c|c|c|}
\hline
One & Two & Three & Four\\
\hline
A & B & C & D\\
E & F & G & H\\
\hline
\end{tabular}
\end{table}

Để tạo bảng biểu, tham khảo thêm tại \url{sharelatex.com}.
Bảng~\ref{tab:table1} là một ví dụ về bảng.
Ngoài ra, có một số tool online~\footnote{ http://www.tablesgenerator.com/} có thể được dùng để tạo bảng biểu một cách trực quan.

\section{Hình vẽ}
Hình \ref{fig_1} là một ví dụ về sử dụng gói graphicx để chèn hình ảnh trong latex. Chú ý là $\backslash${\tt{label}} phải xuất hiện SAU (hoặc cùng) $\backslash${\tt{caption}}.
Khi chèn hình, $\backslash${\tt{caption}} nên xuất hiện sau $\backslash${\tt{includegraphics}}.

\begin{figure}[ht]
\centering
\includegraphics[width=1.5in]{images/logo.jpg}
\caption{Ví dụ về hình ảnh.}
\label{fig_1}
\end{figure}


\subsection{Phương trình toán học}
Việc trình bày phương trình toán học trên một dòng đơn hoặc dòng kép tùy ý, tuy nhiên phải thống nhất trong toàn báo cáo.

Ví dụ về thêm công thức sử dụng ''equation''. Để đánh số các phương trình, hãy sử dụng macro $\backslash${\tt{label}} để tạo một định danh cho phương trình. LaTeX sẽ tự động đánh số phương trình cho bạn.

\begin{equation}
\label{deqn_ex1}
x = \sum_{i=0}^{n} 2{i} Q.
\end{equation}






\chapter{Thuật toán, mã nguồn}
\label{Chapter4}

\section{Thuật toán, mã giả}
Để soạn thảo thuật toán hoặc mã giả trong LaTeX, bạn có thể sử dụng package \textit{algorithm}:
\begin{lstlisting}
\usepackage{algorithm}
\end{lstlisting}

Có thể chèn thuật toán/mã giả như sau:

\begin{algorithm}
\caption{An algorithm with caption}\label{alg:cap}
\begin{algorithmic}[1]
\Require $n \geq 0$
\Ensure $y = x^n$
\State $y \gets 1$
\State $X \gets x$
\State $N \gets n$
\While{$N \neq 0$}
\If{$N$ is even}
    \State $X \gets X \times X$
    \State $N \gets \frac{N}{2}$  \Comment{This is a comment}
\ElsIf{$N$ is odd}
    \State $y \gets y \times X$
    \State $N \gets N - 1$
\EndIf
\EndWhile
\end{algorithmic}
\end{algorithm}

Lưu ý rằng tên lệnh do 'algpseudocode' cung cấp thường được viết hoa chữ cái đầu, ví dụ: $\backslash${\tt{State}}, $\backslash${\tt{While}}, $\backslash${\tt{EndWhile}}.


\section{Mã nguồn}
Để chèn mã nguồn, cần dùng package \textit{listings}:

\begin{lstlisting}
\usepackage{listings}
\end{lstlisting}

Mã nguồn có thể được chèn trực tiếp như sau:

\begin{lstlisting}[language=C++]
cout << "Hello, World!" << endl;
\end{lstlisting}
hoặc chèn thông qua tập tin chứa mã nguồn trong thư mục $SourceCode$ như sau:

\lstinputlisting[language=Python]{SourceCode/hello.py}







\chapter{Trích dẫn tài liệu}
\label{Chapter5}

Mọi ý kiến, khái niệm có ý nghĩa, mang tính chất gợi ý không phải của riêng tác giả và mọi tham khảo khác phải được trích dẫn và chỉ ra nguồn trong danh mục Tài liệu tham khảo của báo cáo. Nguồn được trích dẫn phải được liệt kê chính xác trong danh mục Tài liệu tham khảo.

Việc trích dẫn, tham khảo chủ yếu nhằm thừa nhận nguồn của những ý tưởng có giá trị giúp người đọc theo được mạch suy nghĩ của tác giả, không làm trở ngại việc đọc.

Dùng lệnh $\textbackslash cite$ để trích dẫn một hoặc nhiều tài liệu tham khảo.
Tài liệu tham khảo có thể là trang web~\cite{algorithms2024overleaf}, bài báo khoa học~\cite{le2012human}, sách~\cite{Swanson1999Mathematics}, bài tạp chí~\cite{le2023attention}, bằng sáng chế~\cite{kim2024transformer} hoặc các nguồn tham khảo khác.

% Công trình của tác giả (nếu không có thì comment 02 dòng dưới)
\addcontentsline{toc}{chapter}{Danh mục công trình của tác giả}
\chapter*{Danh mục công trình của tác giả}
\label{Appendix1}

\begin{enumerate}
\item Le, Viet-Tuan and Kim, Yong-Guk. “Attention-based residual autoencoder for video anomaly
detection”. In: Applied Intelligence 53.3 (2023), pp. 3240–3254. issn: 1573-7497. doi: \url{https://doi.org/10.1007/s10489-022-03613-1}.

\end{enumerate}

% In tài liệu tham khảo
\addcontentsline{toc}{chapter}{Tài liệu tham khảo}
\printbibheading[title={Tài liệu tham khảo}]

%\printbibliography[heading=subbibliography, title={Tiếng Việt}, keyword=\textbf{Tieng}, resetnumbers=true]

\DeclareNameAlias{sortname}{last-first}
\DeclareNameAlias{default}{last-first}

\printbibliography[heading=subbibliography, title={Tiếng Anh}, notkeyword=Tieng, resetnumbers=1] 
% ===================================================================== %
% CHÚ Ý: phải gán lại resetnumbers=số tài liệu tham khảo tiếng Việt + 1 %
% ===================================================================== %

% Phần phụ lục
\appendix

\chapter{Tên phụ lục 1}
\label{Appendix1}

Đây là phụ lục.

\end{document} 